% ---
% titre: Triangles semblables
% theme: ES
% chapitre: Figures géométriques planes
% sous_chapitre: Triangles semblables
% niveau: [11e]
% type: EVAL
% tags: [c-triangles-semblables, c-thales, c-calcul-angles,c-perimetre,p-question-TS]
% difficulte: 3
% corrige: true
% ---

\question
Soit la figure ci-dessous où

\begin{itemize}[label=$\triangleright$]
\item $BDEF$ est un trapèze.
\item $AB = 980\,m$; $BC = 490\,m$; $CD = 686\,m$ $BD = 784\,m$
\end{itemize}

\vfill
%schéma
	\begin{center}
	\begin{tikzpicture}[scale=0.9]
	%cercle Omega
	\tkzDefPoints{0/0/O,0/3/Z}
	\tkzDrawCircle[line width=1](O,Z)
	%point A
	\tkzDefPointOnCircle[R = center O angle 180 radius 3cm]]\tkzGetPoint{A}
	\tkzDrawPoints[shape=cross](A)\tkzLabelPoint[left](A){$A$}
	%point B
	\tkzDefPointOnCircle[R = center O angle 0 radius 3cm]]\tkzGetPoint{B}	
	\tkzDrawPoints[shape=cross](B)\tkzLabelPoint[above right](B){$B$}
	%point C
	\tkzDefPointOnCircle[R = center O angle 100 radius 3cm]]\tkzGetPoint{C}
	\tkzDrawPoints[shape=cross](C)\tkzLabelPoint[above](C){$C$}
	%point E
	\tkzDefPointOnCircle[R = center O angle 230 radius 3cm]]\tkzGetPoint{E}	
	\tkzDrawPoints[shape=cross](E)\tkzLabelPoint[below left](E){$E$}
	%point D
	\tkzInterLL(A,B)(C,E)\tkzGetPoint{D}
	\tkzDrawPoints[shape=cross](D)\tkzLabelPoint[above left = 4pt](D){$D$}
	%point F
	\tkzDefLine[parallel=through E](A,B) \tkzGetPoint{X}
	\tkzInterLL(C,B)(E,X)\tkzGetPoint{F}
	\tkzDrawPoints[shape=cross](F)\tkzLabelPoint[below right](F){$F$}
	%Segements
	\tkzDrawSegments(A,B C,F C,E A,E E,F)
	\end{tikzpicture}
	\end{center}

\vfill
\begin{parts}
\part[3]
Démontre que les triangles $CDB$ et $ADE$ sont semblables. Justifie ton raisonnement.

\vspace{10pt}
	{\footnotesize\raggedleft Espace pour ta démarche et tes calculs\par}
	\begin{solutionorgrid}[230pt]
	$\widehat{EAB} = \widehat{ECB}$, car angles inscrits regardant le même arc \hfill \textbf{(1pt)}
	
	\vspace{5pt}
	$\widehat{ADE} = \widehat{BDC}$, car angles opposés par le sommet \hfill \textbf{(1pt)}
	
	\vspace{5pt}
	Les triangles $CDB$ et $ADE$ sont semblables, car deux triangles sont semblables si deux angles sont égaux. \hfill \textbf{(1pt)}
	\end{solutionorgrid}

\begin{tcolorbox}[enhanced jigsaw, interior hidden, colframe=box, parbox=false]%
Réponse: 
\sol{\vspace{20pt}}{Oui, les triangles $CDB$ et $ADE$ sont semblables.}
\end{tcolorbox}	

\newpage
\part[4]
Calcule la longueur de $DE$.

\vspace{10pt}
	{\footnotesize\raggedleft Espace pour ta démarche et tes calculs\par}
	\begin{solutionorgrid}[200pt]
	Comme $CDB$ et $ADE$ sont semblables: 
	
	\vspace{10pt}
	$\dfrac{BC}{AE} = \dfrac{BD}{DE} = \dfrac{CD}{AD}$ \hfill \textbf{(1pt)}
	
	\vspace{10pt}
	$BD = 784\,m$ ; $CD = 686\,m$ ; $AD = 980 - 784 = 196\,m$ \hfill \textbf{(1pt)}%%%%%%%%%%%%%
	
	\vspace{5pt}
	$\dfrac{784}{DE} = \dfrac{686}{196}$ \hfill \textbf{(1pt)}
	
	\vspace{5pt}
	$DE = \dfrac{784 \cdot 196}{686} = 224\,m$ \hfill \textbf{(1pt)}
	\end{solutionorgrid}

\begin{tcolorbox}[enhanced jigsaw, interior hidden, colframe=box, parbox=false]%
Réponse: 
\sol{\vspace{20pt}}{$DE$ vaut 224\,m.}
\end{tcolorbox}	

\vfill
\part[2]
Démontre que les triangles $CDB$ et $CEF$ sont semblables.

\vspace{10pt}
	{\footnotesize\raggedleft Espace pour ta démarche et tes calculs\par}
	\begin{solutionorgrid}[200pt]
	Démontrer que les triangles $CDB$ et $CEF$ sont semblables:
	
	\vspace{10pt}
	$\widehat{DCB} = \widehat{ECF}$, car même angle
	
	\vspace{5pt}
	$\widehat{FEC} = \widehat{BDC}$, car correspondants et $EF // AB$ \hfill \textbf{(1pt)}
	
	\vspace{5pt}
	Les triangles $CDB$ et $CEF$ sont semblables, car deux triangles sont semblables si deux angles sont égaux. \hfill \textbf{(1pt)}
	\end{solutionorgrid}
	
\begin{tcolorbox}[enhanced jigsaw, interior hidden, colframe=box, parbox=false]%
Réponse: 
\sol{\vspace{20pt}}{Oui, les triangles $CDB$ et $CEF$ sont semblables.}
\end{tcolorbox}	

\newpage
\part[6]
Calcule le périmètre du trapèze $BDEF$.

\textit{Si tu n'as pas trouvé la réponse précédente, prend $DE = 200\,m$.}

\vspace{10pt}
	{\footnotesize\raggedleft Espace pour ta démarche et tes calculs\par}
	\begin{solutionorgrid}[500pt]

	Comme $CDB$ et $CEF$ sont semblables: 
	
	\vspace{10pt}
	$\dfrac{CD}{CE} = \dfrac{BC}{CF} = \dfrac{BD}{EF}$ \hfill \textbf{(1pt)}
	
	\vspace{10pt}
	$BD = 784\,m$ ; $CD = 686\,m$ ; $BC = 490\,m$ ; $CE = 686 + 224 = 910\,m$ 
	
	\vspace{5pt}
	$\dfrac{686}{910} = \dfrac{490}{CF} = \dfrac{784}{EF}$ \hfill \textbf{(1pt)}
	
	\vspace{5pt}
	$CF = \dfrac{490 \cdot 910}{686} = 650\,m$ \hfill \textbf{(1pt)}
	
	\vspace{5pt}
	$BF = 650 - 490 = 160\,m$ \hfill \textbf{(1pt)}
	
	\vspace{10pt}
	$EF = \dfrac{784 \cdot 910}{686} = 1040\,m$ \hfill \textbf{(1pt)}
	
	\vspace{20pt}
	Périmètre du trapèze: 
	
	\vspace{5pt}
	$P = BD + DE + EF + BF = 784 + 224 + 1040 + 160 = 2208\,m$ \hfill \textbf{(1pt)}
	\end{solutionorgrid}

\begin{tcolorbox}[enhanced jigsaw, interior hidden, colframe=box, parbox=false]%
Réponse: 
\sol{\vspace{20pt}}{Le périmètre du trapèze vaut 2208\,m.}
\end{tcolorbox}		
\end{parts}